\section{A nonradial channel bound with small total material volume}
\label{nr:section}

The following result controls physical excitation without a bound on the
diameter, location, number of pieces, or minimum feature size of the material
support. It imposes a small \emph{total material volume}, which is an additional
resource absent from Conjecture~\ref{pr:channel-target}. Thus it proves a new
restricted channel theorem, not the unrestricted conjecture. All frequency
estimates in this section are at a positive \emph{real} wavenumber $k$.

Small volume limits the strength of the material's total response, even when
that material is distributed across a large region. The argument first bounds
the free-space response kernel, then uses that bound to control the actual
field in the physical incoming-power norm. Combining this field control with
the earlier flux estimate yields a count. The conclusion is conditional on
the quantitative volume restriction stated below, not on section area alone.

\subsection{An exact real-frequency kernel estimate}

Let $\chi=\varepsilon_r-1$ be supported in a bounded measurable set
$\Omega\subset\mathbb R^3$ of volume $V>0$, with
$0\leq\chi\leq C:=\varepsilon_+-1$. Write
\begin{equation}
 a=\left(\frac{3V}{4\pi}\right)^{1/3},\qquad
 \Phi_k(x)=\frac{e^{ik|x|}}{4\pi|x|},\qquad
 \mathsf K_k=(k^2 I+\nabla\nabla)\Phi_k.
 \label{nr:setup}
\end{equation}
Distributionally on the whole space,
\begin{equation}
 \mathsf K_k=-\mathsf P+\mathsf R_k,\qquad
 \mathsf R_k=k^2\Phi_k I+\nabla\nabla(\Phi_k-\Phi_0),
 \label{nr:split}
\end{equation}
where $\mathsf P$ is the longitudinal orthogonal projection on
$L^2(\mathbb R^3)^3$, with Fourier symbol
$\xi\xi^{\mathsf T}/|\xi|^2$. In particular $\|\mathsf P\|=1$.
The static distributional term, including its diagonal contribution, is
entirely contained in $-\mathsf P$.

\begin{lemma}[PROVED: diameter-independent dynamic kernel estimate]
\label{nr:kernel}
For real $k>0$ and $x\ne0$,
\begin{equation}
 \|\mathsf R_k(x)\|_{\mathbb C^3\to\mathbb C^3}
 \leq\frac{3k^2}{8\pi|x|}.
 \label{nr:kernelbound}
\end{equation}
If $U$ is another bounded measurable set of positive volume and
$b=(3|U|/(4\pi))^{1/3}$, then
\begin{equation}
 \|\mathbf1_U\mathsf R_k\mathbf1_\Omega\|_{2\to2}
 \leq\frac34 k^2ab.
 \label{nr:crossR}
\end{equation}
\end{lemma}
\begin{proof}
Put $r=|x|$, $z=kr$, $n=x/r$, and
$C_z=e^{iz}(1-iz)-1$. Direct radial differentiation gives
\begin{equation}
 \mathsf R_k(x)
 =\frac{z^2e^{iz}-C_z}{4\pi r^3}(I-nn^{\mathsf T})
  +\frac{2C_z}{4\pi r^3}nn^{\mathsf T}.
 \label{nr:radialkernel}
\end{equation}
Since $C_0=0$ and $C_z'=ze^{iz}$, real $z\geq0$ satisfies
$|C_z|\leq z^2/2$. The two real orthogonal projections in
\eqref{nr:radialkernel} give a normal matrix, whose operator norm is the
larger modulus of its two eigenvalues. The tangential modulus is at most
$3z^2/(8\pi r^3)$ and the radial modulus at most $z^2/(4\pi r^3)$.
This proves \eqref{nr:kernelbound}. Near zero the difference
$\Phi_k-\Phi_0=ik/(4\pi)+O(r)$ has a locally integrable Hessian with
no point mass: its first derivative is bounded and its surface integral
on a radius-$r$ sphere is $O(r^2)$. Thus this computation introduces no
unaccounted distributional term.

For any measurable set $A$ of volume $4\pi d^3/3$, and any center $x$,
\begin{equation}
 \int_A |x-y|^{-p}\,dy\leq\frac{4\pi d^{3-p}}{3-p}
 \quad(0<p<3).
 \label{nr:rearrangement}
\end{equation}
Indeed, remove the portion of $A$ outside the radius-$d$ ball about $x$
and fill its equally sized missing portion inside the ball. The decreasing
radial kernel cannot reduce the integral. Integration over the ball gives
the stated expression; truncation justifies the argument at its singularity.
For the kernel majorant in \eqref{nr:kernelbound}, the row integral over
$\Omega$ is at most $3k^2a^2/4$ and the column integral over $U$ at most
$3k^2b^2/4$. The vector-valued Schur estimate, obtained by applying the
scalar Schur argument to the pointwise Euclidean norms, gives
\eqref{nr:crossR}.
\end{proof}

\begin{theorem}[PROVED: real-frequency response controlled by total volume]
\label{nr:volumeinverse}
Suppose
\begin{equation}
 t:=\frac34 C(ka)^2<1.
 \label{nr:smallvolume}
\end{equation}
For a physical incident Maxwell field $e^{\rm inc}$, let $e$ be its
outgoing total electric field in the normalization
$e=\sqrt{\varepsilon_0}E$. Then
\begin{equation}
 \|\sqrt\chi\,e\|_{L^2(\Omega)}
 \leq\frac{\|\sqrt\chi\,e^{\rm inc}\|_{L^2(\Omega)}}{1-t},\qquad
 \|\chi e\|_{L^2(\Omega)}
 \leq\frac{\sqrt C\,\|\sqrt\chi\,e^{\rm inc}\|_{L^2(\Omega)}}{1-t}.
 \label{nr:weightedbound}
\end{equation}
The inverse estimate is uniform over all supports of this volume and all
measurable contrasts in the stated interval. No lower positive bound on
$\chi$ is required.
\end{theorem}
\begin{proof}
The outgoing volume equation is
$e=e^{\rm inc}+\mathsf K_k(\chi e)$.
Set $s=\sqrt\chi$ and $j=se$, with functions on $\Omega$ extended by zero
before applying whole-space operators. Its weighted form is
\begin{equation}
 (I+s\mathsf P s-s\mathsf R_k s)j=se^{\rm inc}.
 \label{nr:weightedVIE}
\end{equation}
The operator $B=I+s\mathsf P s$ is self-adjoint and $B\geq I$, so
$\|B^{-1}\|\leq1$. Lemma~\ref{nr:kernel}, with $U=\Omega$, gives
$\|s\mathsf R_k s\|\leq t$. Factor the left side as
$B(I-B^{-1}s\mathsf R_k s)$ and invert the second factor by its norm
convergent geometric series. This gives an inverse norm at most $(1-t)^{-1}$,
which proves the first inequality; the second follows from
$\|sj\|\leq\sqrt C\|j\|$.

Conversely the solution $j$ reconstructs the outgoing field
$e=e^{\rm inc}+\mathsf K_k(sj)$, and \eqref{nr:weightedVIE} gives $se=j$.
Thus it solves the distributional Maxwell volume equation. A homogeneous
outgoing solution has $j=0$ by the same inverse, and its volume
representation then gives $e=0$. In the piecewise smooth physical class
this is the ordinary unique outgoing scattering solution. The estimate
therefore controls the actual physical field, rather than a formal Born
approximation.
\end{proof}

\subsection{Physical incoming power and the full local field resource}

Take the physical incoming space to be transverse $L^2(S^2)^3$ with
$g(d)\cdot d=0$ and incoming power $\|g\|_2^2$. In the normalized variables
$e=\sqrt{\varepsilon_0}E$, $h=\sqrt{\mu_0}H$,
$c_0=(\varepsilon_0\mu_0)^{-1/2}$, use
\begin{equation}
 \begin{split}
 e^{\rm inc}_g(x)&=\frac{k}{\pi\sqrt{2c_0}}
       \int_{S^2}g(d)e^{ikd\cdot x}\,dS(d),\\
 h^{\rm inc}_g(x)&=\frac{k}{\pi\sqrt{2c_0}}
       \int_{S^2}d\times g(d)e^{ikd\cdot x}\,dS(d).
 \end{split}
 \label{nr:herglotz}
\end{equation}
Stationary phase for smooth $g$ gives incoming electric amplitude of modulus
$\sqrt{2/c_0}|g(-\widehat x)|/r$. Multiplication by the flux factor
$c_0/2$ and integration over the incoming sphere gives $\|g\|_2^2$.
Density extends this power normalization to $L^2$ data. In particular no
plane wave of infinite incoming power is substituted for a unit input.
Cauchy--Schwarz gives, pointwise,
\begin{equation}
 |e_g^{\rm inc}(x)|,\ |h_g^{\rm inc}(x)|
 \leq k\sqrt{\frac{2}{\pi c_0}}\,\|g\|_2.
 \label{nr:incidentbound}
\end{equation}

\begin{theorem}[PROVED: explicit physical-power bound on local electric and magnetic fields]
\label{nr:localresource}
Under \eqref{nr:smallvolume}, fix a bounded measurable observation set $U$
of positive volume, and put
\begin{equation}
 \begin{aligned}
 b&=\left(\frac{3|U|}{4\pi}\right)^{1/3},&
 I_U&=k\sqrt{\frac{2|U|}{\pi c_0}},&
 P_V&=\frac{Ck}{1-t}\sqrt{\frac{2V}{\pi c_0}},\\
 A_k(r)&=r+\frac{k}{2}r^2,&
 D_E&=1+\frac34 k^2ab,&
 D_H&=k\sqrt{A_k(a)A_k(b)}.
 \end{aligned}
 \label{nr:resourceconstants}
\end{equation}
For every physical incoming $g$,
\begin{equation}
 \begin{split}
 \|\chi e_g\|_{L^2(\Omega)}&\leq P_V\|g\|_2,\\
 \|e_g\|_{L^2(U)}&\leq(I_U+D_EP_V)\|g\|_2,\\
 \|h_g\|_{L^2(U)}&\leq(I_U+D_HP_V)\|g\|_2.
 \end{split}
 \label{nr:fieldbounds}
\end{equation}
Consequently the complete local resource satisfies
\begin{equation}
 \int_U(\varepsilon_r|e_g|^2+|h_g|^2)\,dx
 \leq B_U\|g\|_2^2,\qquad
 B_U=\varepsilon_+(I_U+D_EP_V)^2+(I_U+D_HP_V)^2.
 \label{nr:resourcebound}
\end{equation}
These bounds are independent of the distance between $U$ and the material,
and of the diameter and fine structure of the material support.
\end{theorem}
\begin{proof}
Apply \eqref{nr:weightedbound} and \eqref{nr:incidentbound} to obtain the
first inequality. Put $p=\chi e_g$, extended by zero. The exact electric
scattered field is $e^{\rm sc}=-\mathsf Pp+\mathsf R_kp$.
The whole-space projection has norm one, and \eqref{nr:crossR} bounds the
restriction of the dynamic term to $U$. Adding the incident-field bound
gives the second inequality.

Taking the curl of the volume representation and using
$\nabla\times e=ikh$ removes the Hessian part and gives
\begin{equation}
 h^{\rm sc}=\frac{k}{i}\nabla\times(\Phi_k*p),\qquad
 |\nabla\Phi_k(x)|\leq\frac1{4\pi}
                  \left(\frac{k}{|x|}+\frac1{|x|^2}\right).
 \label{nr:magnetickernel}
\end{equation}
The cross-product kernel has operator norm at most $|\nabla\Phi_k|$.
Equation~\eqref{nr:rearrangement}, for $p=1,2$, bounds its row integral
over $\Omega$ by $A_k(a)$ and its column integral over $U$ by $A_k(b)$.
The Schur estimate therefore gives
$\|h^{\rm sc}\|_{L^2(U)}\leq D_H\|p\|_2$.
This proves the magnetic inequality. The coefficient bound
$\varepsilon_r\leq\varepsilon_+$ proves \eqref{nr:resourcebound}.
\end{proof}

\begin{proposition}[PROVED: outgoing scattering is uniformly close to vacuum]
\label{nr:farfield}
With the same normalization and hypotheses, the radiated scattered power
satisfies
\begin{equation}
 P_{\rm sc}(g)\leq
 \left[\frac{Ck^3V}{2\pi(1-t)}\right]^2\|g\|_2^2.
 \label{nr:scatteredpower}
\end{equation}
Equivalently, the full scattering operator differs from the vacuum scattering
operator in power norm by at most $Ck^3V/[2\pi(1-t)]$.
\end{proposition}
\begin{proof}
For $p=\chi e_g$, the normalized electric far field is
\begin{equation}
 e^{\rm sc}_\infty(d)=\frac{k^2}{4\pi}(I-dd^{\mathsf T})
             \int_\Omega p(y)e^{-ikd\cdot y}\,dy.
 \label{nr:farfieldformula}
\end{equation}
Cauchy--Schwarz in $\Omega$ and then integration over $S^2$ give
$P_{\rm sc}=(c_0/2)\|e^{\rm sc}_\infty\|_2^2
\leq c_0k^4V\|p\|_2^2/(8\pi)$.
Insert the first bound in \eqref{nr:fieldbounds} to obtain
\eqref{nr:scatteredpower}. The difference of the outgoing amplitudes is
exactly the scattered outgoing amplitude; the incoming amplitudes agree.
\end{proof}

\subsection{A genuine restricted channel theorem and its remaining limit}

\begin{corollary}[PROVED: uniform channel count under a small-volume cap]
\label{nr:channelcount}
Fix the section, real frequency, material interval, and thresholds of
Definition~\ref{pr:count}, with $\delta<\eta\tau$. Restrict the material
support by $|\Omega|\leq V_{\max}$, where
\begin{equation}
 \frac34(\varepsilon_+-1)k^2
       \left(\frac{3V_{\max}}{4\pi}\right)^{2/3}<1.
 \label{nr:volumecap}
\end{equation}
Use physical incoming amplitudes \eqref{nr:herglotz} and retain the physical
section compatibility of Assumption~\ref{ch:tracecompatibility}. Let $B$
be the fixed comparison box and $K,C_H$ its constants in
\eqref{ch:constants}. Define $B_B$ by \eqref{nr:resourceconstants}--
\eqref{nr:resourcebound}, taking $U=B$ and $V=V_{\max}$. Then every counted
subspace satisfies
\begin{equation}
 \dim W\leq\frac{|B|}{2\pi^2}
      \left(\frac{K C_H B_B}{\eta\tau-\delta}\right)^{3/2}.
 \label{nr:channelbound}
\end{equation}
No enclosing radius, interface count, separation, or minimum feature size
enters this bound.
\end{corollary}
\begin{proof}
All quantities on the right side of \eqref{nr:resourcebound} are
nondecreasing in $V$ while \eqref{nr:smallvolume} holds. Thus the common
volume cap gives the input-to-resource estimate with constant $B_B$.
Apply Theorem~\ref{ch:fulloutput} to any finite-dimensional subspace of $W$.
If $W$ had larger or infinite dimension, it would contain a finite-dimensional
subspace violating the same bound.
\end{proof}

\begin{remark}[OPEN: removal of the volume cap]
\label{nr:scope}
This argument excludes arbitrarily many counted channels formed solely by
redistributing a sufficiently small amount of bounded-contrast dielectric,
even among distant nonradial pieces. It does not exclude large remote
resonators or matching structures: their total material volume is unrestricted
in the primary target. If $t\geq1$, the geometric-series proof supplies no
inverse bound; it does not assert that the scattering problem is singular.
Likewise an $L^2$ bound does not bound pointwise electric peaks and is fully
compatible with the conical hotspot constructions. Finally,
\eqref{nr:kernelbound} uses $|e^{iz}|=1$ for real $z$. At complex resonance
frequencies it gives no diameter-independent estimate, and no resonance-free
complex neighborhood is claimed.
\end{remark}
